\chapter{Tata Kelola Kecerdasan Artifisial}

\section*{Tujuan Pembelajaran}
\begin{itemize}
  \item Mengidentifikasi kasus pemakaian AI (\textit{Artificial Intelligence}) di organisasi dan mengklasifikasi risikonya.
  \item Menjelaskan persyaratan sistem manajemen AI menurut ISO/IEC 42001:2023 dan pendekatan berbasis risiko dalam EU AI Act.
  \item Merumuskan kebijakan pemakaian AI yang dapat dijalankan dan diawasi.
\end{itemize}

\section{Pendahuluan}

% TODO
% Cakupan: staf memakai chatbot AI umum untuk merangkum catatan pasien tanpa izin; organisasi tidak tahu AI apa saja yang sudah dipakai.

\section{Peluang dan Risiko AI di Layanan Kesehatan}

% TODO
% Cakupan: contoh pemakaian AI di rumah sakit dan klinik; jenis risiko (keselamatan pasien, bias, privasi, halusinasi, keamanan); shadow AI; temuan laporan industri 2026 tentang pelanggaran data yang melibatkan AI.

\section{ISO/IEC 42001:2023 dan Standar Pendukung}

% TODO
% Cakupan: struktur sistem manajemen AI dan kontrol Annex A; ISO/IEC 42005:2025 tentang penilaian dampak sistem AI; kaitan dengan ISO/IEC 27001 dan ISO/IEC 27701.

\section{NIST AI RMF}

% TODO
% Cakupan: empat fungsi (Govern, Map, Measure, Manage); profil AI generatif NIST AI 600-1; kapan NIST AI RMF lebih praktis daripada sistem manajemen formal.

\section{EU AI Act dan Pendekatan Berbasis Risiko}

% TODO
% Cakupan: kategori praktik terlarang, risiko tinggi, kewajiban transparansi, dan risiko minimal; kewajiban penyedia dan penerap; jadwal setelah Regulation (EU) 2026/1744 (sistem Annex III mulai 2 Desember 2027, sistem dalam produk Annex I mulai 2 Agustus 2028); relevansinya bagi organisasi Indonesia.

\section{Posisi Regulasi AI di Indonesia}

% TODO
% Cakupan: Surat Edaran Menkominfo 9/2023 tentang etika kecerdasan artifisial; rancangan Perpres tentang peta jalan dan etika AI (status Oktober 2026, periksa ulang); hak keberatan atas keputusan otomatis dalam UU PDP; ketentuan Kemenkes tentang perangkat dan layanan kesehatan berbasis AI (periksa ulang).
% Kasus Medikara: fitur triase gejala dan ringkasan rekam medis dalam kondisi target.

\section{Inventaris AI, Pengawasan Manusia, dan Kebijakan}

% TODO
% Cakupan: inventaris dan klasifikasi kasus pemakaian; bentuk pengawasan oleh manusia; isi kebijakan pemakaian AI dan pelaporan insiden AI; EA untuk AI dibahas rinci di IT40513 dan API untuk agen AI, misalnya Model Context Protocol, di IT40613.
% Pendalaman teknis: menilai kasus berbasis LLM dengan OWASP Top 10 for LLM Applications dan merancang uji evaluasi untuk halusinasi serta kebocoran data.

\section{Aktivitas: Inventaris AI dan Kebijakan Pemakaian AI}

% TODO
% Tujuan aktivitas: mengetahui AI apa saja yang dipakai atau direncanakan, seberapa berisiko, dan aturan pakainya.
% Peran dalam tim: anggota berlatar kesehatan atau hukum menilai dampak pada pasien; analis teknologi menilai data dan model; analis regulasi menautkan kebijakan ke UU PDP; pemilik konteks mendata pemakaian AI yang sudah berjalan.
% Langkah: daftar kasus pemakaian AI yang ada dan yang mungkin muncul di kondisi target; klasifikasi risikonya dengan EU AI Act sebagai pembanding; tetapkan bentuk pengawasan manusia; rumuskan kebijakan (dilarang, perlu persetujuan, pengawasan, pelaporan insiden).
% Luaran dan tonggak capstone: inventaris AI beserta klasifikasi risiko dan kebijakan pemakaian AI; masuk C3.
% Berkas sources yang direncanakan: sources/chapter11/inventaris-ai.xlsx, sources/chapter11/templat-kebijakan-ai.md

\section{Ringkasan}

% TODO

\section{Latihan}

\begin{enumerate}
  \item Klasifikasikan fitur triase gejala Medikara+ menurut kategori risiko EU AI Act. Jelaskan alasannya dan kewajiban yang muncul.
  \item Jelaskan mengapa larangan total memakai AI generatif sering justru memperbanyak \textit{shadow AI}.
  \item Bandingkan ISO/IEC 42001 dan ISO/IEC 27001 dari sisi tujuan dan jenis risiko yang dikelola.
\end{enumerate}
